\documentclass[UTF8,aspectratio=169]{ctexbeamer}
\usetheme{Madrid}
\usecolortheme{default}
\usepackage{amsmath}
\usepackage{booktabs}
\usepackage{siunitx}
\usepackage{graphicx}

\title{Engineered-Feature MLP for Spray Penetration\\ (v2: Hard Physics Prior on Pressure and Diameter)}
\author{Mie Spray Postprocessing Project}
\date{\today}

\begin{document}

\begin{frame}
  \titlepage
\end{frame}

\begin{frame}{Message}
\textbf{Claim:} the sparse-feature instability shown for the v1 direct-feature MLP can be removed at its source by replacing the pressure/diameter inputs with a single hard physical scaling group taken from the Hiroyasu--Arai / Zhou penetration laws, and by training the network on the resulting non-dimensional target.

\vspace{0.4em}
\textbf{What this deck does:}
\begin{itemize}
  \item recap the v1 sparse-feature failure mode
  \item state the physical scaling group used in v2
  \item describe exactly which features are removed, which are added, and how the target is reparameterized
  \item describe the new training scaffolding (group splits, pretrain collapse check, two variants)
\end{itemize}
\end{frame}

\begin{frame}{Recap: What v1 Showed}
\textbf{v1 used 9 direct features:}
\[
[\text{time}, \theta_{\mathrm{tilt}}, n_{\mathrm{plumes}}, d_n, t_i, \log P_{\mathrm{inj}}, \log P_{\mathrm{ch}}, \log\Delta P, P_{\mathrm{cb}}]
\]
\textbf{Observed pathology (see \texttt{slides\_sparse\_feature\_instability}):}
\begin{itemize}
  \item nozzle diameter has only 6 levels and is not uniformly available across pressures
  \item dense MLP sweeps over $d_n$ produce $5$--$6$ gradient sign changes and $p95\,|d^2\mu/dd_n^2|\sim10^7$
  \item injection pressure shows worst-secant-deviation ratio $>1$ on weakly supported slices
\end{itemize}

\vspace{0.3em}
\textbf{Diagnosis:} the network is allowed to invent continuous interpolation between sparse levels because the data does not constrain $\partial\mu/\partial d_n$ or $\partial\mu/\partial\log P_{\mathrm{inj}}$ almost anywhere.
\end{frame}

\begin{frame}{Engineering Move: Use the Theory Directly}
The Hiroyasu--Arai post-breakup branch and the Zhou et al.\ EOI extension share one well-supported scaling group. Reading the post-breakup branch
\[
S(t)\;\propto\;K_p\!\left(\frac{\Delta P}{\rho_a}\right)^{1/4}d_n^{1/2}\,t^{1/2},
\]
the entire dependence on injection pressure, chamber pressure, chamber density and nozzle diameter collapses to one scalar
\[
\boxed{\;A\;=\;\left(\frac{\Delta P}{\rho_a}\right)^{1/4}d_n^{1/2}\;}
\]
which has units of length and acts as a per-condition length scale.

\vspace{0.3em}
\textbf{Idea:} stop asking the network to learn $d_n$, $\log P_{\mathrm{inj}}$, $\log P_{\mathrm{ch}}$, $\log\Delta P$ as four separate sparse inputs. Use $A$ as a hard divisor on the target and feed the network the non-dimensional trajectory.
\end{frame}

\begin{frame}{Computing $A$ from the Test Matrix}
\textbf{Per curve, $A$ is determined entirely from the operating point:}
\begin{itemize}
  \item $d_n$ from \texttt{nozzle\_properties.diameter\_mm} in the test matrix JSON
  \item $\Delta P=P_{\mathrm{inj}}-P_{\mathrm{ch}}$ from the per-row pressures
  \item $\rho_a$ from the linear law $\rho_a=\SI{5}{kg/m^3}\times P_{\mathrm{ch}}/\SI{4.42}{bar}$, calibrated against the explicit \texttt{chamber\_density\_kg\_per\_m3} entries provided for Nozzle1--8
\end{itemize}

\vspace{0.4em}
\textbf{Implementation reference:}
\begin{itemize}
  \item \texttt{engineered\_feature\_common.build\_canonical\_feature\_table} computes \texttt{ambient\_pressure\_bar\_phys}, \texttt{ambient\_density\_kg\_m3}, \texttt{delta\_pressure\_bar\_phys}, \texttt{A\_scale}, \texttt{log\_A}
  \item the canonicalization is the same for Nozzle0 (derived) and Nozzle1--8 (looked up), so all datasets share one consistent $A$
\end{itemize}
\end{frame}

\begin{frame}{New Feature Set}
\textbf{Removed (4 sparse direct features):}
\[
\{\,d_n,\;\log P_{\mathrm{inj}},\;\log P_{\mathrm{ch}},\;\log\Delta P\,\}
\]

\textbf{Two variants are exposed in \texttt{FEATURE\_COLUMNS\_BY\_VARIANT}:}
\begin{itemize}
  \item \texttt{a\_only} (5 features):
  \[
  [\text{time}_{\mathrm{norm}},\;\theta_{\mathrm{tilt},z},\;n_{\mathrm{plumes},z},\;t_{i,z},\;P_{\mathrm{cb},z}]
  \]
  $A$ never enters as an input; it only acts as a divisor on the target.
  \item \texttt{a\_plus\_log\_a} (6 features):
  \[
  [\text{time}_{\mathrm{norm}},\;\theta_{\mathrm{tilt},z},\;n_{\mathrm{plumes},z},\;t_{i,z},\;P_{\mathrm{cb},z},\;\log A_z]
  \]
  $A$ also appears once as a one-dimensional residual axis.
\end{itemize}

\textbf{Why two variants:} \texttt{a\_only} is the strongest commitment to the Zhou scaling. \texttt{a\_plus\_log\_a} lets the network express residual dependence on the scaling group as a single normalized axis instead of four sparse axes.
\end{frame}

\begin{frame}{Target Reparameterization}
\textbf{v1:} the network predicted $\mu(t)$ and $\log\sigma^2(t)$ on the physical penetration $S(t)$ in mm.

\vspace{0.3em}
\textbf{v2:} the network predicts the same two heads on the non-dimensional target
\[
\hat S(t)\;=\;S(t)/A,\qquad
\hat\sigma^2(t)\;=\;\sigma^2(t)/A^{2}.
\]
At inference, the physical mean and standard deviation are recovered exactly by
\[
\mu_S(t)\;=\;A\,\hat\mu(t),\qquad
\sigma_S(t)\;=\;A\,\hat\sigma(t).
\]

\vspace{0.3em}
\textbf{Implementation reference:} \texttt{PenetrationDataset} stores \texttt{target\_scaled = penetration / A\_scale} per row. The Stage-1 MSE and Stage-2 NLL objectives in \texttt{engineered\_feature\_common} consume \texttt{target\_scaled} directly, while \texttt{toy\_inference\_from\_run} multiplies the predicted mean and std by \texttt{A\_scale} when reporting in millimetres.
\end{frame}

\begin{frame}{Why This Should Help the Sparse-Feature Pathology}
\textbf{Mechanism:} the four sparse axes that produced spurious curvature in v1 are no longer free inputs. Their entire effect is now expressed through one analytic factor $A$ whose exponents are fixed by theory rather than being learned from a $6$-level grid.

\vspace{0.4em}
\textbf{Expected consequences:}
\begin{itemize}
  \item $\partial\mu_S/\partial d_n$ inherits the closed-form $d_n^{1/2}$ shape; the network has no remaining degrees of freedom to invent oscillations along $d_n$
  \item $\partial\mu_S/\partial P_{\mathrm{inj}}$, $\partial\mu_S/\partial P_{\mathrm{ch}}$ inherit the closed-form $(\Delta P/\rho_a)^{1/4}$ shape, with positive sensitivity by construction
  \item all remaining learning capacity is spent on the dense, well-supported axes (time, tilt, plumes, injection duration, control backpressure)
\end{itemize}

\vspace{0.3em}
\textbf{Honest caveat:} this only helps if the Zhou scaling actually collapses the trajectories on the present data. That is why a pretrain collapse check is now part of the training driver.
\end{frame}

\begin{frame}{Pretrain Collapse Check}
Before any optimization, \texttt{run\_pretrain\_collapse\_check} reconstructs the representative trajectories from the four fitted constants, divides each by its own $A$, and tests whether
\[
\hat S(t)\;\approx\;\text{function of the remaining (dense) features only}.
\]

\vspace{0.3em}
\textbf{What it produces:}
\begin{itemize}
  \item per-time-slice variance of $\hat S(t)$ versus the original $S(t)$
  \item a linear regression $R^{2}$ of $\hat S(t)$ on the kept features
  \item a binary \texttt{passed} flag, surfaced in the run directory
\end{itemize}

\vspace{0.3em}
\textbf{Role in the pipeline:} if the precheck fails, training aborts unless \texttt{--allow-failed-precheck} is passed. This makes the physical assumption \emph{falsifiable} on the actual dataset rather than asserted.
\end{frame}

\begin{frame}{Other Pipeline Changes}
\textbf{Group-aware splits.} \texttt{assign\_splits\_by\_group} assigns whole \texttt{sample\_group\_id} groups to train/val/test instead of permuting individual rows. This prevents leakage between curves that share an operating point.

\vspace{0.3em}
\textbf{Stage-1 / Stage-2 driver split.} \texttt{train\_stage1\_mse.py} and \texttt{train\_stage2\_nll.py} replace the v1 monolithic notebooks. Stage 2 warm-starts from a Stage-1 engineered run directory and reuses its scaler state, so the variant choice and the $A$-scaling are guaranteed identical between stages.

\vspace{0.3em}
\textbf{Loss objectives.} \texttt{compute\_loss\_stage1\_engineered} and \texttt{compute\_loss\_stage2\_engineered\_nll} keep the v1 monotonicity and post-transient concavity priors, but apply them on $\hat\mu(t)$ instead of $\mu(t)$. The shape constraints transfer cleanly because dividing by a positive constant $A$ preserves sign and curvature.

\vspace{0.3em}
\textbf{Inference path.} \texttt{toy\_inference\_from\_run.py} now branches on \texttt{infer\_feature\_family}: legacy v1 runs follow the 9-feature path, engineered v2 runs canonicalize the raw input through the same $\rho_a$, $\Delta P$, $A$ pipeline used in training and rescale the model outputs by $A$ before plotting.
\end{frame}

\begin{frame}{Engineering Takeaway}
\textbf{What v2 commits to:}
\begin{itemize}
  \item the Hiroyasu--Arai post-breakup scaling group $A=(\Delta P/\rho_a)^{1/4}d_n^{1/2}$ is treated as a hard prior, not a soft regularizer
  \item the network is no longer asked to learn a continuous law over a $6$-point diameter grid or a $4$-point pressure grid
  \item the resulting model is falsifiable: if the collapse check fails, the prior is wrong on this dataset and training stops
\end{itemize}

\vspace{0.4em}
\textbf{What still has to be evaluated post-training:}
\begin{itemize}
  \item rerun the v1 sparse-feature instability diagnostics (\texttt{compare\_sparse\_feature\_stability}) and confirm that the secant-deviation ratios and curvature norms drop materially
  \item verify that the held-out NLL on physical $S$ is at least as good as v1 once the Jacobian factor $\log A$ is accounted for
  \item compare \texttt{a\_only} versus \texttt{a\_plus\_log\_a} on diagnostics, not just on validation loss
\end{itemize}
\end{frame}

\end{document}